\documentclass[border=4pt]{standalone}
\usepackage[siunitx]{circuitikz}

\begin{document}
\begin{circuitikz}[european resistors]
  \draw (0,0) node[ground] {}
    to[V, l_=$V_S$] (0,3)
    -- (2,3)
    to[R, l=$R_1$, -*] (2,1.5)
    to[R, l=$R_2$] (2,0)
    -- (0,0);

  % Source terminal test point; the dot marks its three-conductor junction.
  \node[circ] at (1,3) {};
  \draw (1,3) -- (1,3.75)
    node[ocirc, label={[font=\small]above:{TP\_S}}] {};

  % The external load is genuinely in parallel with R_2.
  \draw (2,1.5) -- (5,1.5)
    to[R, l=$R_L$] (5,0)
    -- (2,0);
  \node[circ] at (2,0) {};

  % Output polarity, placed away from symbols and conductors.
  \node[font=\small] at (4.58,1.25) {$+$};
  \node[font=\small] at (4.58,0.25) {$-$};
  \node[font=\small] at (4.02,0.75) {$V_{\mathrm{out}}$};

  % Test points brought out to clear space.
  \draw (2,1.5) -- (3,2.35)
    node[ocirc, label={[font=\small]above:{TP1}}] {};
  \draw (2,0) -- (3,-0.85)
    node[ocirc, label={[font=\small]below:{TP0 ($0$~V)}}] {};
\end{circuitikz}
\end{document}
